%%%%%
%
\section{Temp�rature de l'air � partir de la temp�rature sonique et
de l'humidit� absolue}

\itemstyle{Op\'erateur~:} {\sf Tair\_soni\_ha}
\index{Tair\_soni\_ha}

%%%%%OBJET%%%%
\itemstyle{Objet~:}
Cet op�rateur calcule la temp�rature statique de l'air � partir de 
la temp�rature sonique et de l'humidit� absolue. 

%%%%%ARGUMENTS%%%%
\bigskip
%%%%%ENTREES%%%%
\begin{tabarg}
\multicolumn{3}{l}{\titretab {Entr\'ees}}\\
\hline
$T_{soni}$ & Vecteur, {\tt float} & Temp�rature sonique ($K$ ou $^\circ C$)\\
$H_a$ & Vecteur, {\tt float} & Humidit� absolue ($g/m^3$)\\
$P$ & Vecteur, {\tt float} & Pression  ($hPa$)\\
\end{tabarg}

\medskip
%%%%%SORTIES%%%%
\begin{tabarg}
\multicolumn{3}{l}{\titretab {Sorties}}\\
\hline
$T$ & Vecteur, {\tt float} & Temp�rature de l'air (m\^eme unit� que $T_{soni}$)\\
\end{tabarg}


%%%%%FORMULE%%%%
\itemstyle{Formule}

%%Tair_soni_ha
Etant donn� ces 3 formules thermodynamiques :
\begin{align*}
T &= (1 - 0.518.q).T_{soni}\\
q &= r/(1+r) \\
r &= H_a.R_a.T/P 
\end{align*}

o� $q$ est l'humidit� sp�cifique, $r$ le rapport de m�lange en vapeur d'eau,
$R_a$ est la constante sp�cifique de l'air sec.\\

On en d�duit que la temp�rature de l'air $T$ est une racine du polyn\^ome du 
second degr� suivant :
\[
T^2 + \left [ \frac{P}{H_a.R_a} - (1 - 0.518)\times T_{soni} \right ].T
    - \frac{P}{H_a.R_a}.T_{soni} = 0  
\]



%%%%%COMMENTAIRE%%%%
\itemstyle{Commentaire~:}
 On appelle temp�rature "sonique" une temp�rature calcul�e
� partir de la vitesse de propagation du son dans l'air
en assimilant l'air � un gaz parfait sec : $c = \sqrt{ \gamma_a R_a T}$
donne $T = \frac{c^2}{\gamma_a R_a}$.

%%%%%AUTEUR%%%%
\itemstyle{Auteur~:}
CNRM/GMEI/TRAMM

%%%%%REFERENCE%%%%
\itemstyle{R\'ef\'erences~:}

Triplet-Roche \cite{Triplet-Roche},
Pr\'esentation \`a l'atelier intrumentation de 1996 \cite{AIE96-TAA}.
